\scanpage{39}
\mspage{35}
% archive sheet number: 38
% source lines 1--8
Il reste la \unsure{perplexité}, sur la terminologie parallèle strate -- multistrates. Finalement, je vais appeler les $X\in M$ les \uline{strates} de l'atelier $M$. Chaque strate détermine l'ensemble $\widetilde X$ de ses \uline{sous-strates}, ou \uline{strates incidentes}, C'est $\widetilde X$ qui sera la \uline{multistrate associée à \unsure{les|la}} \uline{strate} $X$, ou déterminée, définie par $X$, \deleted{\unsure{ou la}}
% source lines 9--16
Ainsi, pour une figure $F$, on distingue entre les \uline{strates de} $F$, qui sont les éléments de $\widetilde F$, ou \uline{strates} de $M$ \uline{incidentes} à $F$, et les \uline{multistrates de} $F$, qui sont les $\widetilde X$ pour $X\in\widetilde F$. On pourrait désigner l'ens de ces multistrates par $\widetilde{\widetilde F}$.

% source lines 17--19
Quand on décrit l'atelier des figures \unsure{élémentaires} ensemblistes sur un ensemble $\unsure{L|\mathcal L}$
\[
\text{Fig \unsure{él|el}}(L)=M
\]
% source lines 20--26
il est entendu que ses objets apparaissent non comme des parties, mais \inserted{disons} comme \deleted{des} ensembles de parties de $L$, \deleted{soit $\widetilde X$}. Cela n'empêche qu'un tel ensemble de parties de $L$, considéré comme objet $\in M$, sera appelé ``strate'', et non ``multistrate''. Ainsi,
